\documentclass[10pt,a4paper]{article}

\usepackage[utf8]{inputenc}
\usepackage[margin=0.7in]{geometry}
\usepackage[T1]{fontenc}
\usepackage{lmodern}
\usepackage{hyperref}
\usepackage{enumitem}
\usepackage{titlesec}
\usepackage{xcolor}

% Color definitions for clean, modern academic aesthetics
\definecolor{primary}{RGB}{20, 40, 75}       % Deep navy for headings
\definecolor{linkcolor}{RGB}{25, 95, 170}     % Professional academic blue for hyperlinks
\definecolor{tagcolor}{RGB}{95, 95, 95}       % Muted gray for brackets and separators

\hypersetup{
    colorlinks=true,
    linkcolor=linkcolor,
    citecolor=linkcolor,
    filecolor=linkcolor,
    urlcolor=linkcolor,
    pdfauthor={Milton Llera Montero, PhD},
    pdftitle={Curriculum Vitae - Milton Llera Montero}
}

% Section styling
\titleformat{\section}
  {\large\bfseries\color{primary}}
  {}{0em}
  {}[\vspace{-0.5em}\rule{\textwidth}{0.75pt}\vspace{-0.3em}]

\titleformat{\subsection}
  {\normalsize\bfseries\color{primary}}
  {}{0em}
  {}

\titlespacing*{\section}{0pt}{1.1em}{0.55em}
\titlespacing*{\subsection}{0pt}{0.75em}{0.35em}

\setlength{\parindent}{0pt}
\setlength{\parskip}{0.35em}

% Custom command for clean publication link badges
\newcommand{\linkbadge}[2]{%
  \textcolor{tagcolor}{[}\href{#2}{\textbf{#1}}\textcolor{tagcolor}{]}%
}

\begin{document}

% ==================== HEADER ====================
\begin{center}
    {\LARGE \textbf{\color{primary} Milton Llera Montero, PhD}} \\[0.35em]
    Currently independent \ $\cdot$ \ 
    \href{mailto:milton.llera@icloud.com}{milton.llera@icloud.com} \ $\cdot$ \ 
    \href{mailto:mlle@itu.dk}{mlle@itu.dk} \\[0.25em]
    \href{mlmontero.com}{Personal Site} \ $\cdot$ \
    \href{https://scholar.google.com/citations?user=mTcxP70AAAAJ&hl=en}{Google Scholar} \ $\cdot$ \ 
    \href{https://www.github.com/miltonllera}{GitHub} \ $\cdot$ \ 
    \href{https://www.linkedin.com/in/milton-l-montero/}{LinkedIn} \ $\cdot$ \
    \href{https://orcid.org/0000-0002-7440-3510}{ORCID} \ $\cdot$ \
    \href{https://www.youtube.com/@miltonllera1437/featured}{YouTube}
\end{center}

\vspace{-0.2em}

% ==================== RESEARCH SUMMARY ====================
\section*{Summary}
I am broadly interested in the mechanisms that support generalization and adaptability in both biological and artificial systems. Specifically, my research has covered a broad range of topics including Cognitive Science, Computer Vision, Artificial Life and Synthetic Biology. Recently, I was a Postdoctoral Researcher at the REAL Lab with Sebastian Risi studying theoretical models of self-organization in the context of biological development. Previously I completed my PhD in Psychology at the University of Bristol under the supervision of Jeffrey Bowers, Rui Ponte Costa, and Gaurav Malhotra, where I studied the combinatorial generalization properties and (dis)similarities between humans and Deep Learning models of vision. My work combines computational modeling, mainly using Deep Learning, along with behavioral and mechanistic interpretability approaches to answer fundamental questions in the aforementioned domains of inquiry.

% ==================== EXPERIENCE ====================
\section*{Experience}

\textbf{Postdoctoral Researcher} \hfill \textit{June 2023 -- September 2026} \\
REAL Lab, IT University of Copenhagen \hfill Copenhagen, Denmark \\
Studying the mechanisms that enable robust, open-ended development and adaptation in biological and artificial systems.

\vspace{0.45em}
\textbf{Applied Scientist Intern} \hfill \textit{July 2022 -- December 2022} \\
Causal Representation Learning Lablet, Amazon Web Services (AWS) \hfill T\"ubingen, Germany \\
Worked as an intern focusing on action recognition from videos using motion cues with Carl-Johann Simon-Gabriel, Thomas Brox, and Francesco Locatello.

\vspace{0.45em}
\textbf{Teaching Assistant} \hfill \textit{October 2018 -- July 2020} \\
Department of Computer Science, University of Bristol \hfill Bristol, UK \\
Teaching assistant for \textit{Information Processing in the Brain} (COMSM034), covering information theory, deep learning, and data analysis relevant to computational neuroscience.

\vspace{0.45em}
\textbf{Research Assistant} \hfill \textit{August 2017 -- May 2018} \\
School of Engineering, Technical University of Madrid (UPM) \hfill Madrid, Spain \\
Implemented software tools for the Human Brain Project through the Cajal Blue Brain affiliate project.

% ==================== EDUCATION ====================
\section*{Education}

\textbf{Doctor of Philosophy (PhD) in Psychology} \hfill \textit{September 2018 -- May 2023} \\
School of Psychological Science, University of Bristol \hfill Bristol, UK \\
Supervisors: Prof. Jeffrey S. Bowers, Dr. Rui Ponte Costa, Dr. Gaurav Malhotra \\
\textit{Thesis:} \linkbadge{Combinatorial Generalisation in Machine Vision}{https://research-information.bris.ac.uk/en/studentTheses/combinatorial-generalisation-in-machine-vision/}

\vspace{0.45em}
\textbf{Master's Degree (MSc) in Artificial Intelligence} \hfill \textit{September 2016 -- July 2017} \\
School of Engineering, Technical University of Madrid (UPM) \hfill Madrid, Spain \\
Supervisors: Prof. Concha Bielza, Prof. Pedro Larrañaga \hfill \textit{CGPA: 8.8/10} \\
\textit{Thesis:} \linkbadge{A Novel Multi-dimensional Regression Model Based on Gaussian Networks}{https://cig.fi.upm.es/wp-content/uploads/2024/01/TFM_MILTON_LLERA_MONTERO.pdf}

\vspace{0.45em}
\textbf{Bachelor's Degree (BS) in Computer Science} \hfill \textit{September 2017 -- July 2016} \\
Department of Mathematics and Computer Science, University of Havana \hfill Havana, Cuba \\
Supervisor: Dr. Alejandro Piad Morffis \hfill \textit{CGPA: 4.52/5} \\
\textit{Thesis:} Example-based Terrain Modeling for Procedural Content Generation.

% ==================== PUBLICATIONS ====================
\section*{Publications}

\subsection*{Journal Articles}
\begin{enumerate}[leftmargin=1.6em, itemsep=0.45em]
    \item Bowers, J.~S., Malhotra, G., Dujmovi\'{c}, M., \textbf{Montero, M.~L.}, Tsvetkov, C., Biscione, V., \dots, \& Blything, R. (2023). \textbf{Deep problems with neural network models of human vision}. \textit{Behavioral and Brain Sciences}, 46, e385. \\
    \linkbadge{Target Article DOI}{https://doi.org/10.1017/S0140525X22002813}

    \item Bowers, J.~S., Malhotra, G., Dujmovi\'{c}, M., \textbf{Montero, M.~L.}, Tsvetkov, C., Biscione, V., \dots, \& Blything, R. (2023). \textbf{Clarifying status of DNNs as models of human vision} (Authors' Response). \textit{Behavioral and Brain Sciences}, 46, e412. \\
    \linkbadge{Response DOI}{https://doi.org/10.1017/S0140525X23002777}

    \item Bowers, J.~S., Malhotra, G., Adolfi, F., Dujmovi\'{c}, M., \textbf{Montero, M.~L.}, Biscione, V., \dots, \& Heaton, R.~F. (2023). \textbf{On the importance of severely testing deep learning models of cognition}. \textit{Cognitive Systems Research}, 82, 101158. \\
    \linkbadge{Journal DOI}{https://doi.org/10.1016/j.cogsys.2023.101158}

    \item Bykowska, O., Gontier, C., Sax, A.-L., Jia, D.~W., \textbf{Montero, M.~L.}, Bird, A.~D., Houghton, C., Pfister, J.-P., \& Costa, R.~P. (2019). \textbf{Model-based inference of synaptic transmission}. \textit{Frontiers in Synaptic Neuroscience}, 11, 21. \\
    \linkbadge{Journal DOI}{https://doi.org/10.3389/fnsyn.2019.00021}

    \item \textbf{Montero, M.~L.}, Sacramento, J., \& Costa, R.~P. (2019). \textbf{Computational roles of plastic probabilistic synapses}. \textit{Current Opinion in Neurobiology}, 54, 90--97. \\
    \linkbadge{Journal DOI}{https://doi.org/10.1016/j.conb.2018.09.002}
\end{enumerate}

\subsection*{Peer-Reviewed Conference Proceedings}
\begin{enumerate}[leftmargin=1.6em, itemsep=0.45em]
    \item \textbf{Montero, M.~L.}, Najarro, E., Schauser, J., \& Risi, S. (2026). \textbf{Learning developmental scaffolds for self-organising systems}. In \textit{Artificial Life Conference Proceedings 38} (Vol. 2026, No. 1, p. 36). MIT Press. \\
    \linkbadge{MIT Press Proceedings}{https://direct.mit.edu/isal/proceedings/isal2026/38/44/138194} \ $\cdot$ \ 
    \linkbadge{arXiv:2605.14998}{https://arxiv.org/abs/2605.14998}

    \item Nisioti, E., Plantec, E., \textbf{Montero, M.}, Pedersen, J., \& Risi, S. (2025). \textbf{When Does Neuroevolution Outcompete Reinforcement Learning in Transfer Learning Tasks?}. In \textit{Proceedings of the Genetic and Evolutionary Computation Conference (GECCO 2025)}, pp. 48--57. ACM. \\
    \linkbadge{ACM Proceedings (PDF)}{https://dl.acm.org/doi/pdf/10.1145/3712256.3726475}

    \item Pedersen, J.~W., Plantec, E., Nisioti, E., Barylli, M., \textbf{Montero, M.~L.}, Korte, K., \& Risi, S. (2025). \textbf{Hypernetworks That Evolve Themselves}. In \textit{Artificial Life Conference Proceedings 37 (ALife 2025)}, Vol. 2025, No. 1, p. 17. MIT Press. \\
    \linkbadge{MIT Press Proceedings}{https://direct.mit.edu/isal/proceedings/isal2025/37/17/134104} \ $\cdot$ \ 
    \linkbadge{arXiv:2512.16406}{https://arxiv.org/abs/2512.16406}

    \item \textbf{Montero, M.~L.}, Plantec, E., Nisioti, E., Pedersen, J.~W., \& Risi, S. (2024). \textbf{Meta-learning an evolvable developmental encoding}. In \textit{Artificial Life Conference Proceedings 36} (Vol. 2024, No. 1, p. 50). MIT Press. \\
    \linkbadge{MIT Press Proceedings}{https://direct.mit.edu/isal/proceedings/isal2024/36/50/123543} \ $\cdot$ \ 
    \linkbadge{Conference Program}{https://2024.alife.org/detailed_program.html}

    \item Plantec, E., Pedersen, J.~W., \textbf{Montero, M.~L.}, Nisioti, E., \& Risi, S. (2024). \textbf{Evolving Self-Assembling Neural Networks: From Spontaneous Activity to Experience-Dependent Learning}. In \textit{Artificial Life Conference Proceedings 36} (Vol. 2024, No. 1, p. 37). MIT Press. \\
    \linkbadge{MIT Press Proceedings}{https://direct.mit.edu/isal/proceedings/isal2024/36/37/123555}

    \item Pedersen, J.~W., Plantec, E., Nisioti, E., \textbf{Montero, M.~L.}, \& Risi, S. (2024). \textbf{Structurally flexible neural networks: Evolving the building blocks for general agents}. In \textit{Proceedings of the Genetic and Evolutionary Computation Conference (GECCO 2024)}, pp. 1119--1127. ACM. \\
    \linkbadge{ACM Proceedings (PDF)}{https://dl.acm.org/doi/pdf/10.1145/3638529.3654113} \ $\cdot$ \ 
    \linkbadge{arXiv:2404.15193}{https://arxiv.org/abs/2404.15193}

    \item Nisioti, E., Plantec, E., \textbf{Montero, M.~L.}, Pedersen, J., \& Risi, S. (2024). \textbf{Growing artificial neural networks for control: the role of neuronal diversity}. In \textit{Proceedings of the Genetic and Evolutionary Computation Conference Companion (GECCO 2024 Companion)}, pp. 175--178. ACM. \\
    \linkbadge{ACM Proceedings (PDF)}{https://dl.acm.org/doi/pdf/10.1145/3638530.3654356}

    \item \textbf{Montero, M.~L.}, Bowers, J.~S., Ludwig, C.~J., Costa, R.~P., \& Malhotra, G. (2022). \textbf{Lost in Latent Space: Disentangled Models and the Challenge of Combinatorial Generalisation}. In \textit{Proceedings of the 36th Conference on Neural Information Processing Systems (NeurIPS 2022)}. New Orleans, USA. \\
    \linkbadge{NeurIPS Program \& Poster}{https://neurips.cc/virtual/2022/poster/52961} \ $\cdot$ \ 
    \linkbadge{OpenReview}{https://openreview.net/pdf?id=7yUxTNWyQGf} \ $\cdot$ \ 
    \linkbadge{arXiv:2204.02283}{https://arxiv.org/abs/2204.02283}

    \item \textbf{Montero, M.~L.}, Ludwig, C.~J., Costa, R.~P., Malhotra, G., \& Bowers, J.~S. (2020). \textbf{The role of disentanglement in generalisation}. In \textit{International Conference on Learning Representations (ICLR 2021)}. \\
    \linkbadge{ICLR Program \& Poster}{https://iclr.cc/virtual/2021/poster/3173} \ $\cdot$ \ 
    \linkbadge{OpenReview}{https://openreview.net/pdf?id=qbH974jKUVy}
\end{enumerate}

\subsection*{Preprints \& Workshop Papers}
\begin{enumerate}[leftmargin=1.6em, itemsep=0.45em]
    \item Hartl, B., \textbf{Montero, M.~L.}, Barilly, M., Risi, S., \& Levin, M. (2026). \textbf{On Growth and Form, and Function: Reusable Regulatory Handles Control Phenotypic Variation}. \textit{Preprint}. \\
    \linkbadge{arXiv:2609.29755}{https://arxiv.org/pdf/2609.29755}

    \item \textbf{Montero, M.~L.}, Bowers, J., \& Malhotra, G. (2024). \textbf{Successes and Limitations of Object-centric Models at Compositional Generalisation}. In \textit{NeurIPS 2024 Workshop on Compositional Learning: Perspectives, Methods, and Paths Forward}. \\
    \linkbadge{arXiv:2412.18743}{https://arxiv.org/pdf/2412.18743} \ $\cdot$ \ 
    \linkbadge{Workshop Talk (SlidesLive)}{https://slideslive.com/39031181/successes-and-limitations-of-objectcentric-models-at-compositional-generalisation?ref=folder-163258}

    \item Biscione, V., \textbf{Montero, M.~L.}, Dujmovi\'{c}, M., Yin, D., Malhotra, G., Puebla, G., \dots, \& Bowers, J.~S. (2024). \textbf{MindSet: Vision. A toolbox for testing DNNs on key psychological experiments}. \textit{Preprint}. \\
    \linkbadge{arXiv:2404.05290}{https://arxiv.org/pdf/2404.05290}
\end{enumerate}

% ==================== TALKS & PRESENTATIONS ====================
\section*{Talks \& Presentations}

\subsection*{Panel Discussions \& Invited Talks}
\begin{itemize}[leftmargin=1.6em, itemsep=0.35em]
    \item \textbf{No, NeuroAI is not adopting the right methods to advance our understanding of the mind and brain}. GAC Kickoff debate presentation, \textit{Cognitive Computational Neuroscience (CCN 2026)}, New York, USA. \\
    \linkbadge{Talk \& Slides}{https://youtu.be/kYESqAKLGhA?si=vNGPbsdhHVvX91J8&t=131} \ $\cdot$ \ 
    \linkbadge{Panel Discussion}{https://youtu.be/kYESqAKLGhA?si=8GGyCruh40tr0-QB&t=2514}

    \item \textbf{Morphogenetic Scaffolds in Self-organizing Systems}. Invited talk at the Levin Lab, March 2026. (Online). \\
    \linkbadge{Talk \& Slides (YouTube)}{https://youtu.be/X9dK7UKBICE?si=zwaakCn40DeKZKt0}

    \item \textbf{Panel Discussion on Compositional Learning}. Panel speaker alongside Thomas Kipf (Google DeepMind) and Yilun Du (Harvard University). \textit{Workshop on Compositional Learning: Perspectives, Methods, and Paths Forward, NeurIPS 2024}, Vancouver, Canada. \\
    \linkbadge{Panel Discussion (SlidesLive)}{https://slideslive.com/39031182/panel-discussion?ref=folder-163258}

    % \item \textbf{Meta-Learning an Evolvable Developmental Encoding}. Plenary session talk at the \textit{International Conference on Artificial Life (ALife 2024)}, Copenhagen, Denmark, 2024.
\end{itemize}

\subsection*{Selected Contributed Presentations \& Posters}
\begin{itemize}[leftmargin=1.6em, itemsep=0.35em]
    \item \textbf{Language Cellular Automata: Guiding Morphogenesis with Natural Language}. Poster presentation, \textit{ALife Proceedings Companion} (pp. 99--101), Kyoto, Japan, 2025. \\
    \linkbadge{Conference Program}{https://drive.google.com/file/d/1Sg_MLUVagck_DdcrU3eNWUGTpK8RsYRT/view}

    \item \textbf{Scaffolding Development via Morphogenetic Gradients}. Poster presentation, \textit{ALife Proceedings Companion} (pp. 287--289), Kyoto, Japan, 2025. \\
    \linkbadge{Conference Program}{https://drive.google.com/file/d/1Sg_MLUVagck_DdcrU3eNWUGTpK8RsYRT/view}

    \item Vassallo, G., Nisioti, E., Pedersen, J.~W., Plantec, E., \textbf{Montero, M.~L.}, \& Risi, S. (2025). \textbf{Curriculum learning with a hierarchical cellular encoding}. In \textit{Proceedings of the Genetic and Evolutionary Computation Conference Companion (GECCO 2025 Companion)}. ACM. \\
    \linkbadge{ACM Proceedings (PDF)}{https://dl.acm.org/doi/pdf/10.1145/3712255.3726672} 

    \item \textbf{Division versus subtraction: task-specific gating in recurrent neural networks}. \textit{Cognitive and Systems Neuroscience (Cosyne 2021)}, Online, 2021. \\
    \linkbadge{Conference Program}{https://static1.squarespace.com/static/6102ca347474c263c40150cd/t/610870a432c9d257a80cace4/1627943077187/Cosyne2021_program_book.pdf}  \ $\cdot$ \
    \linkbadge{Talk \& Slides (YouTube)}{https://www.youtube.com/watch?v=r8kVb1yDV44}

    \item \textbf{Deep Generative Models as Perceptual Front-ends for Decision Making}. Poster presentation, \textit{53rd Annual Meeting of the Society for Mathematical Psychology (MathPsych 2020)}, Online. \\
    \linkbadge{YouTube Video}{https://youtu.be/9Sje1N35dM0?si=ciZ4ExxLT_0BC-UO}

    \item \textbf{Subtractive gating improves generalization in working memory tasks}. \textit{Conference on Cognitive Computational Neuroscience (CCN 2019)}, Berlin, Germany, 2019. \\
    \linkbadge{Extended Abstract}{https://pdfs.semanticscholar.org/fb2b/329a6d27bcea8f58749f2ff04c353402ad8d.pdf}
\end{itemize}

% ==================== WORKSHOPS ORGANIZED ====================
\section*{Workshops and Special Sessions Organized}
\begin{itemize}[leftmargin=1.6em, itemsep=0.35em]
    \item \textbf{Artificial Life, Intelligence, Complexity \& Evolution (ALICE)}. Geilo, Norway, 2027. \\
    \linkbadge{Workshop Website}{https://aliceworkshop.org}

    \item \textbf{ALife for Science and Engineering}. Special Session at \textit{ALife 2026}, Waterloo, Canada. \\
    \linkbadge{Special Session Website}{https://alifeforscience.github.io}

    \item \textbf{Generative Adversarial Collaboration: Is NeuroAI adopting the right methods and theoretical frameworks to advance our understanding of mind and brain?}. \textit{CCN 2026}, New York, USA. \\
    \linkbadge{GAC Website}{https://sites.google.com/ccneuro.org/gac2020/gacs-by-year/2026-gacs/2026-1?authuser=0}

    \item \textbf{Artificial Life, Intelligence, Complexity \& Evolution (ALICE)}. Copenhagen, Denmark, 2026. \\
    \linkbadge{Workshop Website}{https://2026.aliceworkshop.org}

    \item \textbf{Self-Organizing Neural Intelligences}. Special Session at \textit{ALife 2025}, Kyoto, Japan. \\
    \linkbadge{Special Session Website}{https://sites.google.com/view/soni-alife}

    \item \textbf{Evolving Self-Organisation}. Workshop at \textit{GECCO 2025}, Malaga, Spain. \\
    \linkbadge{Workshop Website}{https://gecco-2025.sigevo.org/Workshopb2fc.html?itemId=5097}
\end{itemize}

% ==================== SKILLS & SERVICE ====================
\section*{Skills \& Academic Service}

\textbf{Technical Skills:} Extensive experience with programming in Python (primary language). Within the AI ecosystem, proficient with both PyTorch and JAX, as well as libraries for scientific computing (NumPy, SciPy) and visualization (Matplotlib). Familiarity with C/C++, C\#, R, JavaScript, and MATLAB.

\vspace{0.35em}
\textbf{Behavioral Experiments:} Design and execution of calibrated psychophysics experiments using PsychoPy; experience with pre-registration frameworks such as the Open Science Framework (OSF) model.

\vspace{0.35em}
\textbf{Scientific Engagement \& Reviewing:} Active reviewer for NeurIPS, ICLR, ICML, AAAI, GECCO, ALife, and TMLR. Experienced in academic speaking and public dissemination.

\vspace{0.35em}
\textbf{Language Proficiency:} Native speaker of both English and Spanish.

\end{document}
